\section{Alfred T1S Hardware Platform Strategy}
\label{sec:platform-strategy}

The E2B endpoints and the Gateway do not use identical silicon, and one PCB does not become both by population. The platform is a shared \emph{electrical and mechanical philosophy}, realized as sibling node classes.

\begin{diagram}
                  ALFRED T1S PLATFORM

             common network philosophy
             common power philosophy
             common connectors
             common protection
             common mechanical language
             common instrumentation
             common validation
             common manufacturing/test approach
                         |
          +--------------+---------------+
          |              |               |
          v              v               v

     E2B-CONTROL     E2B-LIGHTING      GATEWAY

       LAN8660         LAN8661       MCU + LAN8651
          |               |              |
      actuator          lighting     CAN/CAN-FD/LIN
\end{diagram}

\subsection{What is shared, and what is not}

\begin{center}
\small
\begin{tabularx}{\linewidth}{@{}p{3.6cm} X@{}}
\toprule
\textbf{Shared where technically sensible} & \textbf{Deliberately not shared} \\
\midrule
12/24\,V input architecture & Network silicon (LAN8660 vs.\ LAN8661 vs.\ LAN8650/1) -- role-specific by design \\
Protection components/philosophy (reverse-polarity, TVS, overcurrent) & MCU presence -- E2B-Control and E2B-Lighting have none; Gateway requires one \\
T1S connector and termination strategy & Peripheral-side interface circuitry (motor-driver path vs.\ LED-driver path vs.\ CAN/CAN-FD/LIN transceivers) \\
Board outline and mounting pattern, where practical & Firmware -- E2B endpoints run none at all; only the Gateway has an Alfred-authored firmware image \\
Status-indication style & \\
Current/voltage measurement approach and accessible test points & \\
Labeling conventions & \\
Fabrication process and test-fixture philosophy & \\
Documentation structure & \\
\bottomrule
\end{tabularx}
\end{center}

\begin{decision}[title={Explicit guardrail}]
Do not introduce additional MCU silicon into the LAN8660/LAN8661 E2B designs merely to make them resemble the Gateway board. Any argument for adding a companion MCU to an E2B node must be justified by an actual E2B requirement (Sections~\ref{sec:lan8660-validation}, \ref{sec:lan8661-validation}) discovered during Rev~A/B bring-up, never by a desire for platform symmetry.
\end{decision}

\subsection{Why this is still ``one platform''}

A customer-facing board from any of the three node classes should be visually and electrically recognizable as one product family (shared connector, shared power input behavior, shared status-LED language, shared test-point layout convention) even though the PCBs are not derived from one shared Gerber file. The commonality that actually matters commercially -- one 10BASE-T1S electrical layer, one validation discipline, one manufacturing/test philosophy -- survives intact; the commonality that does not matter (identical silicon) is not manufactured artificially.
